\section{Free-monoid indistinguishability}
\label{sec:clone-theorem}

Let
\begin{equation}
 \MZ=(\N_{>0},\mid,1,\lcmop)
 \label{eq:integer-monoid}
\end{equation}
and let \(P=\At(\MZ)\), the covers of \(1\).  Write the formal free
commutative monoid additively as
\begin{equation}
 \Ffree=\bigoplus_{p\in P}\N e_p,
 \label{eq:formal-monoid}
\end{equation}
with coordinatewise order, zero bottom, and coordinatewise maximum as join.
Define
\begin{equation}
 \PhiVal(n)=(v_p(n))_{p\in P}.
 \label{eq:valuation-map}
\end{equation}
Equation~\eqref{eq:valuation-map} is used to prove an isomorphism; it is not
an operation made available to a candidate kernel.

Transport every admitted decoration:
\begin{equation}
 F'_N=\PhiVal(F_N),\qquad
 w'(\PhiVal(n))=w(n),\qquad
 G'_{\PhiVal(m),\PhiVal(n)}=G_{m,n},
 \label{eq:transport}
\end{equation}
and apply the same index transport to each compiled tensor or operator
coefficient.

\begin{theorem}[Decorated free-monoid indistinguishability]
\label{thm:clone}
The valuation map \(\PhiVal\) is an isomorphism in \(\Csrc\) from the
decorated integer-divisibility source to the transported formal
free-commutative/UFD source.  Consequently, every isomorphism-natural
invariant \(I\) of the admitted data satisfies
\begin{equation}
 I(\MZ)=\PhiVal^*I(\Ffree).
 \label{eq:invariant-equality}
\end{equation}
The conclusion holds for local or nonlocal \(I\), for every arity, and for
scalar-, kernel-, ledger-, operator-coefficient-, or function-valued output.
\end{theorem}
\begin{proof}
Unique factorization makes \(\PhiVal\) a bijection between positive integers
and finitely supported exponent vectors.  Moreover,
\begin{equation}
 m\mid n
 \quad\Longleftrightarrow\quad
 v_p(m)\le v_p(n)\quad\text{for every }p,
 \label{eq:order-transport}
\end{equation}
so it preserves and reflects the order and sends \(1\) to the zero vector.
Covers of \(1\) go exactly to coordinate unit vectors.  For every finite
family,
\begin{equation}
 v_p(\lcmop(n_1,\ldots,n_k))=\max_i v_p(n_i),
 \label{eq:join-transport}
\end{equation}
so all finite joins and all pointed intervals are preserved.  Interval
M{\"o}bius functions therefore agree.  Equation~\eqref{eq:transport} makes
the cutoff and analytic decorations commute with the map by construction.
Thus \(\PhiVal\) is an isomorphism of the complete displayed object.
Naturality of \(I\) at this isomorphism gives
\eqref{eq:invariant-equality}.  No step restricts the locality, arity, or
algebraic form of \(I\).
\end{proof}

\begin{corollary}[The universal selectivity specification is inconsistent]
\label{cor:selectivity}
There is no admissible natural invariant \(I\) such that
\begin{equation}
 I(\MZ)\ne0
 \quad\text{and}\quad
 I(U)=0\quad\text{for every free-commutative/UFD control }U.
 \label{eq:selectivity-gate}
\end{equation}
\end{corollary}
\begin{proof}
The universal quantifier in \eqref{eq:selectivity-gate} includes the
transported clone \(\Ffree\).  The second requirement sets its value to
zero, while \cref{thm:clone} identifies that value with the nonzero baseline
value.
\end{proof}

The clone is not a tuned control.  Comparing an object with a transported
copy is the minimal coherence test for an invariant advertised as natural.
If a control policy forbids transported coefficient data, it no longer
quantifies over every decorated UFD object.  That is a weaker and different
claim.
